\pagenumbering{arabic}

\chapter{Introducción} \label{cap1}
%\addcontentsline{toc}{chapter}{Introducción}

\vspace*{5mm}

%Justificación -> auge IA%
Vivimos una época de auge en el campo de la Inteligencia Artificial provocado por la llegada de nuevas tecnologías en modelos de lenguaje y sistemas generativos que han creado un aumento en la demanda y la inversión lo que ha ocasionado que se desarrollen aún más éstas tecnologías, se haya normalizado su uso y se estén integrando en muchos y variados ámbitos tanto cotidianos (agentes de voz y de texto, asistentes, buscadores...) como en los negocios (asistentes a la programación, automatizaciones, sistemas integrados en medicina e investigación, logística, hostelería...).

%necesidad de nuevas arquitecturas y regulación%
Este florecimiento ha traído consigo también nuevos desafíos tanto en arquitectura como en regulaciones. Que cada vez haya más dispositivos recolectando y produciendo datos necesarios para entrenar modelos hace que estos sistemas tengan que adaptarse para ser más flexibles y poder dar cobertura a diferentes tipos de topologías, por ejemplo, varios dispositivos móviles pueden colaborar para entrenar una red teniendo diferentes especificaciones. Asimismo la compartición de datos personales e información sensible trae consigo la necesidad de crear sistemas más seguros, por ejemplo una persona que esté entrenando un modelo con su geolocalización pero sin tener que compartir su posición con una empresa externa. Este tipo de problemas son el germen para el desarrollo de sistemas de aprendizaje federados. La idea detrás de esta arquitectura es la de que varios clientes entrenen modelos locales y compartan los parámetros (por ejemplo, los pesos de la red) con un servidor de manera que todos cooperan en el entrenamiento conjunto de la red mientras que los datos nunca abandonan los servidores.

%modelos federados heterogéneos%
La bondad de los modelos federados está probada para datos homogéneos

%Beneficios modelos heterogéneos, facilidad integración con modelos ya existentes%
%Problema de los modelos heterogéneos%







%ejemplo cita ~\cite{20NGDataset:web}

El resto de la memoria se organiza del siguiente modo:

\begin{description}
\item[2. Estado del arte.] En este capítulo se analizarán ...
\item[3. Propuesta.] Se describen los sistemas desarrollados en este trabajo. 
\item[4. Estudio Experimental.] Se presentan los resultados de los experimentos realizados....
\item[5. Conclusiones y Trabajo Futuro.] Se plantean las conclusiones y futuras ampliaciones del sistema y líneas de trabajo futuro en este campo.
\end{description}


%%
\cleardoublepage